\chapter{Freie Halbgruppen und allgemeine Darstellungen}

\section{Wörter, freie Halbgruppen und Präsentationen}

\begin{definition}[Freie Halbgruppe und Präsentation]
Für eine Menge $A$ bezeichnet $A^+$ die Menge aller nichtleeren
endlichen Wörter über $A$, versehen mit der Verkettung.
$A^*=A^+\cup\{\varepsilon\}$ enthält zusätzlich das leere Wort,
das nur als äußerer Kontext zugelassen wird.
Für $A=\varnothing$ ist $A^+=\varnothing$.
Zu einer Relation $R\subseteq A^+\times A^+$ heißt
\[
 \langle A\mid R\rangle=A^+/\operatorname{Cg}_{A^+}(R)
\]
die \emph{durch Erzeuger $A$ und Relationen $R$ präsentierte Halbgruppe}.
Weder $A$ noch $R$ müssen endlich sein.
\end{definition}

\begin{theorem}[Universelle Eigenschaft der freien Halbgruppe]
\label{b28guide:thm:free}
$A^+$ ist eine Halbgruppe.
Jede Abbildung $u:A\to S$ besitzt genau eine Fortsetzung
zu einem Homomorphismus $\widehat u:A^+\to S$, nämlich
\[
 \widehat u(a_1\cdots a_n)=u(a_1)\cdots u(a_n).
\]
Sein Bild ist $\langle u(A)\rangle$.
Jede andere Halbgruppe mit derselben universellen Eigenschaft
ist durch einen eindeutig bestimmten, die Erzeuger respektierenden
Isomorphismus zu $A^+$ isomorph.
\cite{harju,gould}
\end{theorem}

\begin{theorem}[Jede Halbgruppe besitzt eine Präsentation]
\label{b28guide:thm:presentation}
Für jede Halbgruppe $S$ und jedes Erzeugendensystem $A\subseteq S$
ist die Auswertung $\pi:A^+\to S$ surjektiv, und
\[
 S\cong A^+/\ker\pi .
\]
Insbesondere liefert ein Alphabet
$\{\underline{s}:s\in S\}$ mit den Relationen
\[
 \underline{s}\,\underline{t}=\underline{st}
 \qquad(s,t\in S)
\]
eine Präsentation von $S$.
\cite{harju}
\end{theorem}

\begin{theorem}[Homomorphismen aus einer Präsentation]
Eine Abbildung $u:A\to T$ setzt sich genau dann zu einem
Homomorphismus $\langle A\mid R\rangle\to T$ fort, wenn
\[
 \widehat u(v)=\widehat u(w)\qquad\text{für alle }(v,w)\in R.
\]
Die Fortsetzung ist dann eindeutig.
Zwei Wörter repräsentieren genau dann dasselbe Element der
präsentierten Halbgruppe, wenn sie durch eine endliche Kette
von Ersetzungen aus $R$ oder $R^{-1}$ in Wortkontexten
$x\,\square\,y$ mit $x,y\in A^*$ verbunden sind.
\cite{harju}
\end{theorem}

\begin{theorem}[Genau die durch Assoziativität erzwungenen Identitäten]
\label{b28guide:thm:identities}
Seien $t,u$ zwei Terme aus Variablen und einer zweistelligen
Verknüpfung, ohne Konstanten.
Die Gleichung $t=u$ gilt bei jeder Belegung in jeder Halbgruppe
genau dann, wenn nach dem Entfernen aller Klammern in beiden
Termen dieselbe Folge von Variablen in derselben Reihenfolge steht.
Damit ist jede allgemeingültige Identität dieser Sprache bereits
durch Assoziativität bestimmt.
\cite{harju,burris}
\end{theorem}

\section{Freie Produkte beliebiger Halbgruppen}

\begin{definition}[Freies Halbgruppenprodukt]
Für zwei Halbgruppen $S,T$ betrachten wir disjunkte Kopien ihrer Mengen.
Die Menge $S*T$ besteht aus allen nichtleeren endlichen Folgen von
Elementen dieser Kopien, in denen benachbarte Glieder verschiedenen
Faktoren entstammen. Das Produkt wird durch Aneinanderhängen gebildet;
entstammen die beiden Glieder an der Verbindungsstelle demselben Faktor,
werden sie durch ihr Produkt in diesem Faktor ersetzt.
\end{definition}

\begin{theorem}[Normalform und universelle Eigenschaft des freien Produkts]
\label{b28guide:thm:free-product}
$S*T$ ist eine Halbgruppe. Ihre Elemente besitzen genau eine
alternierende Normalform wie in der Definition.
Die kanonischen Abbildungen $i_S:S\to S*T$ und $i_T:T\to S*T$
sind Einbettungen. Für jede Halbgruppe $U$ und Homomorphismen
$f:S\to U$, $g:T\to U$ existiert genau ein Homomorphismus
$h:S*T\to U$ mit
\[
 h\circ i_S=f,\qquad h\circ i_T=g.
\]
Einheiten einzelner Faktoren werden weder identifiziert noch
aus den Normalformen gestrichen. Es gilt auch
$\varnothing*T\cong T$ und $S*\varnothing\cong S$.
\cite{nybergbrodda2023,howie}
\end{theorem}
\section{Universelles amalgamiertes Produkt}

\begin{definition}[Pushout von Halbgruppenhomomorphismen]
Zu Homomorphismen $\alpha:U\to S$ und $\beta:U\to T$ heißt
eine Halbgruppe $P$ mit Homomorphismen $j_S:S\to P$ und $j_T:T\to P$
ein \emph{Pushout}, wenn $j_S\circ\alpha=j_T\circ\beta$ gilt
und folgende universelle Eigenschaft erfüllt ist:
Zu jeder Halbgruppe $V$ und Homomorphismen $f:S\to V$, $g:T\to V$
mit $f\circ\alpha=g\circ\beta$ existiert genau ein Homomorphismus
$h:P\to V$ mit
\[
 f=h\circ j_S,\qquad g=h\circ j_T.
\]
Der Pushout heißt auch \emph{universelles amalgamiertes Produkt}.
Alle Abbildungen müssen nur die Multiplikation erhalten;
die Erhaltung einer vorhandenen Einheit wird nicht verlangt.
\end{definition}

\begin{theorem}[Existenz universeller amalgamierter Produkte]
\label{b28guide:thm:pushout}
Für beliebige Halbgruppen $U,S,T$ und Homomorphismen
$\alpha:U\to S$, $\beta:U\to T$ existiert ein Pushout.
Mit den kanonischen Einbettungen $i_S:S\to S*T$ und $i_T:T\to S*T$
ist er gegeben durch
\[
\begin{aligned}
 \rho&=\operatorname{Cg}_{S*T}
 \bigl(\{(i_S(\alpha(u)),i_T(\beta(u))):u\in U\}\bigr),\\
 P&=(S*T)/\rho,\qquad
 j_S=q_\rho\circ i_S,\quad j_T=q_\rho\circ i_T.
\end{aligned}
\]
Er ist bis auf den eindeutig bestimmten, die Strukturabbildungen
respektierenden Isomorphismus eindeutig.
Die Abbildungen $j_S,j_T$ müssen nicht injektiv sein,
selbst wenn $\alpha,\beta$ Einbettungen sind.
Für $U=\varnothing$ ergibt sich das freie Produkt $S*T$.
\cite{howie1968}
\end{theorem}

\section{Treue Darstellungen durch Abbildungen}

\begin{definition}[Transformationshalbgruppe und Wirkung]
Für eine Menge $X$ ist $\mathcal T_X$ die Menge aller, auch
nichtbijektiver, Abbildungen $X\to X$ mit der üblichen Komposition
$(f\circ g)(x)=f(g(x))$.
Eine \emph{linke Wirkung} von $S$ auf $X$ ist eine Abbildung
$S\times X\to X$, $(a,x)\mapsto a\cdot x$, mit
$(ab)\cdot x=a\cdot(b\cdot x)$.
Ein zusätzliches Einheitswirkungsgesetz wird nicht verlangt.
Sie heißt \emph{treu}, wenn verschiedene Elemente von $S$
verschiedene Abbildungen auf $X$ bewirken.
\end{definition}

\begin{theorem}[Wirkungen und Transformationsdarstellungen]
$\mathcal T_X$ ist eine Halbgruppe.
Die linken Wirkungen von $S$ auf $X$ entsprechen genau den
Homomorphismen $S\to\mathcal T_X$ durch
$a\mapsto(x\mapsto a\cdot x)$.
Treue Wirkungen entsprechen Einbettungen.
\cite{gould,howie}
\end{theorem}

\begin{theorem}[Cayleys Darstellungssatz für Halbgruppen]
\label{b28guide:thm:cayley}
Jede Halbgruppe $S$ besitzt eine treue Transformationsdarstellung.
Konkret ist
\[
 \lambda:S\longrightarrow\mathcal T_{S^1},\qquad
 a\longmapsto\lambda_a,\qquad \lambda_a(x)=ax,
\]
eine Einbettung; es gilt $\lambda_{ab}=\lambda_a\circ\lambda_b$.
Somit ist jede Halbgruppe zu einer Unterhalbgruppe einer
Transformationshalbgruppe isomorph.
\cite{harju,gould}
\end{theorem}

\section{Einbettung in eine einfache Halbgruppe}

\begin{definition}[Bruck-Erweiterung]
Für eine beliebige Halbgruppe $S$ setzen wir
$B(S)=\mathbb N_0\times S^1\times\mathbb N_0$,
wobei $\mathbb N_0=\{0,1,2,\ldots\}$, mit der Verknüpfung
\[
 (i,a,j)(k,b,\ell)=
 \begin{cases}
 (i-j+k,b,\ell),&j<k,\\
 (i,ab,\ell),&j=k,\\
 (i,a,\ell+j-k),&j>k.
 \end{cases}
\]
\end{definition}

\begin{theorem}[Einbettungssatz von Bruck]
\label{b28guide:thm:bruck}
Für jede Halbgruppe $S$ ist $B(S)$ eine einfache Halbgruppe
mit neutralem Element $(0,1,0)$. Die Abbildung
\[
 S\longrightarrow B(S),\qquad s\longmapsto(0,s,0)
\]
ist eine Einbettung. Somit läßt sich jede Halbgruppe in eine
einfache Halbgruppe mit Einheit einbetten.
Einfachheit und Einheit betreffen die konstruierte Erweiterung;
an $S$ wird keine dieser Bedingungen gestellt.
\cite{munn1968}
\end{theorem}
\clearpage
\section{Weitere allgemeine Einbettungssätze}

\begin{definition}[Eigenschaften einer Erweiterung]
Für eine Halbgruppe $T$ sei $E(T)=\{e\in T:e^2=e\}$.
Sie heißt \emph{idempotentenerzeugt}, wenn $T=\langle E(T)\rangle$,
und \emph{bisimpel}, wenn sie genau eine $\mathcal D$-Klasse besitzt.
Sie heißt \emph{kongruenzfrei}, wenn
$\operatorname{Con}(T)=\{\Delta_T,T\times T\}$ gilt.
Diese Begriffe beschreiben die Zielhalbgruppen der folgenden
Einbettungen. Die Einbettungen erhalten die Multiplikation;
eine vorhandene Einheit von $S$ muß keine Einheit von $T$ werden.
\end{definition}

\begin{theorem}[Einbettungssatz von Higgins]
\label{b28guide:thm:higgins}
Für jede Halbgruppe $S$ existieren eine Halbgruppe $T$ und
eine Einbettung $i:S\to T$ mit
\[
 T=E(T)E(T).
\]
Jedes Element von $T$ ist somit ein Produkt zweier idempotenter Elemente
von $T$. Die Faktoren dürfen außerhalb von $i(S)$ liegen;
Idempotente oder solche Produktdarstellungen innerhalb von $S$
werden nicht vorausgesetzt.
\cite{higgins1995,oliveira2014}
\end{theorem}

\begin{theorem}[Einbettungssatz von Howie]
\label{b28guide:thm:howie-embedding}
Jede Halbgruppe $S$ besitzt eine Einbettung in eine bisimple,
idempotentenerzeugte, kongruenzfreie Halbgruppe $T$.
Gleichzeitig gelten also
\[
 |T/\mathcal D|=1,\qquad
 T=\langle E(T)\rangle,\qquad
 \operatorname{Con}(T)=\{\Delta_T,T\times T\}.
\]
Alle drei Eigenschaften gehören zur Erweiterung $T$;
ihre idempotenten Erzeuger dürfen außerhalb des Bildes von $S$ liegen.
An $S$ wird keine zusätzliche algebraische Eigenschaft
oder Beschränkung seiner Mächtigkeit gestellt.
\cite{howie1981}
\end{theorem}

\section{Unterscheidung durch Kontexte}

\begin{definition}[Syntaktische Kongruenz einer Teilmenge]
Für eine beliebige Teilmenge $P\subseteq S$ sei
\[
 a\equiv_P b
 \iff
 (\forall x,y\in S^1)\,
 \bigl(xay\in P\iff xby\in P\bigr).
\]
Eine Kongruenz \emph{sättigt} $P$, wenn $P$ eine Vereinigung
ihrer Äquivalenzklassen ist.
Der Quotient $S/{\equiv_P}$ heißt \emph{syntaktische Halbgruppe}
der Teilmenge $P$.
\end{definition}

\begin{theorem}[Größte Kongruenz mit vorgegebener Sättigung]
\label{b28guide:thm:syntactic}
$\equiv_P$ ist die größte Kongruenz auf $S$, die $P$ sättigt.
Erkennt ein Homomorphismus $f:S\to T$ die Menge $P$, das heißt
$P=f^{-1}(Q)$ für ein $Q\subseteq T$, so gilt
\[
 \ker f\subseteq{\equiv_P}.
\]
Es gibt genau einen surjektiven Homomorphismus
$f(S)\to S/{\equiv_P}$, der $f(a)$ auf $[a]_{\equiv_P}$ abbildet.
\cite{harju96}
\end{theorem}

\clearpage
\section{Subdirekte Zerlegung}

\begin{definition}[Subdirekte Darstellung und Irreduzibilität]
Eine Einbettung $h:S\to\prod_{i\in I}T_i$ heißt \emph{subdirekt},
wenn jede Abbildung $\pi_i\circ h:S\to T_i$ surjektiv ist.
Eine Halbgruppe $T$ mit mindestens zwei Elementen heißt
\emph{subdirekt irreduzibel}, wenn in jeder subdirekten Darstellung
von $T$ mindestens eine Koordinatenabbildung injektiv ist.
\end{definition}

\begin{theorem}[Kongruenzkriterium für subdirekte Darstellungen]
Für jede Familie $(\rho_i)_{i\in I}$ von Kongruenzen ist
\[
 h:S\to\prod_{i\in I}S/\rho_i,\qquad
 a\mapsto([a]_{\rho_i})_{i\in I}
\]
ein Homomorphismus mit surjektiven Koordinatenabbildungen und
$\ker h=\bigcap_i\rho_i$.
Er ist genau dann eine subdirekte Darstellung, wenn
$\bigcap_i\rho_i=\Delta_S$.
Eine Halbgruppe mit mindestens zwei Elementen ist genau dann
subdirekt irreduzibel, wenn sie eine kleinste Kongruenz besitzt,
die echt größer als $\Delta_S$ ist.
\cite{burris,howie}
\end{theorem}

\begin{theorem}[Subdirekter Darstellungssatz von Birkhoff]
\label{b28guide:thm:subdirect}
Jede Halbgruppe ist zu einem subdirekten Produkt homomorpher Bilder
isomorph, von denen jedes subdirekt irreduzibel ist oder höchstens
ein Element besitzt. Für eine Halbgruppe mit mindestens zwei
Elementen können alle Faktoren subdirekt irreduzibel gewählt werden.
Die leere und die einelementige Halbgruppe können jeweils selbst
als einziger Faktor dienen.

Dieser Existenzsatz wird im üblichen mengentheoretischen Rahmen
mit Auswahlaxiom formuliert; die Zerlegung muß weder endlich noch
eindeutig sein.
\cite{burris}
\end{theorem}

\clearpage
\section{Die feinste Halbverbandszerlegung}

\begin{definition}[Halbverbandskongruenz und Unzerlegbarkeit]
Ein \emph{Halbverband} ist eine kommutative idempotente Halbgruppe.
Eine Kongruenz $\rho$ auf $S$ heißt \emph{Halbverbandskongruenz},
wenn $S/\rho$ ein Halbverband ist.
Eine nichtleere Halbgruppe heißt \emph{halbverbandsunzerlegbar},
wenn jedes ihrer homomorphen Bilder, das ein Halbverband ist,
genau ein Element besitzt.
\end{definition}

\begin{theorem}[Größter Halbverbandsquotient]
\label{b28guide:thm:semilattice-quotient}
Für jede Halbgruppe $S$ ist
\[
 \eta=\operatorname{Cg}_S\bigl(
 \{(a^2,a):a\in S\}\cup
 \{(ab,ba):a,b\in S\}\bigr)
\]
die kleinste Halbverbandskongruenz.
Jeder Homomorphismus $f:S\to Y$ in einen Halbverband $Y$
faktorisiert eindeutig als $f=\bar f\circ q_\eta$ durch
einen Homomorphismus $\bar f:S/\eta\to Y$.
Insbesondere ist jeder Halbverbandsquotient von $S$
ein homomorphes Bild von $S/\eta$.
\cite{bogdanovic2011}
\end{theorem}

\begin{theorem}[Zerlegungssatz von Tamura]
\label{b28guide:thm:tamura}
Für jede Halbgruppe $S$ ist jede Klasse der Kongruenz $\eta$
eine halbverbandsunzerlegbare Unterhalbgruppe.
Mit $Y=S/\eta$ und $S_\alpha=q_\eta^{-1}(\{\alpha\})$ gilt
\[
 S=\mathop{\dot\bigcup}_{\alpha\in Y}S_\alpha,\qquad
 S_\alpha S_\beta\subseteq S_{\alpha\beta}.
\]
Dies ist die \emph{feinste Halbverbandszerlegung}: Jede andere
Halbverbandskongruenz faßt ausschließlich ganze $\eta$-Klassen zusammen.
Für $S=\varnothing$ sind $Y$ und die Zerlegung leer.
Die Kommutativität und Idempotenz gelten für den Quotienten $Y$;
sie werden weder für $S$ noch für die Komponenten $S_\alpha$ verlangt.
\cite{bogdanovic2011}
\end{theorem}